% Рисунок из тетради: компакт K строго внутри Omega, d = dist(K, dOmega).
\begin{figure}[htbp]
\centering
\begin{tikzpicture}[scale=1.0]
  \draw[thick, fill=blue!4]
    plot[smooth cycle, tension=0.8]
      coordinates {(0,0) (2.4,0.8) (4.4,0.2) (5.0,-1.4) (3.2,-2.6) (0.9,-2.2) (-0.7,-1.0)};
  \node at (4.3,-0.6) {$\Omega$};

  \filldraw[pattern=north east lines, pattern color=blue!45, draw=blue!60!black, thick]
    plot[smooth cycle, tension=0.9]
      coordinates {(1.3,-0.5) (2.6,-0.2) (3.3,-1.0) (2.5,-1.8) (1.4,-1.6)};
  \node at (2.3,-1.0) {$K$};

  \draw[-{Latex[length=2mm]}] (3.3,-1.0) -- node[above, font=\scriptsize]{$d$} (4.75,-1.05);
  \node[anchor=north west, font=\footnotesize] at (3.1,-2.3) {$d=\dist(K,\partial\Omega)>0$};

  \draw[dashed, gray]
    plot[smooth cycle, tension=0.9]
      coordinates {(1.15,-0.3) (2.65,0.0) (3.55,-1.0) (2.6,-2.0) (1.25,-1.75)};
\end{tikzpicture}
\caption{Компакт $K\subset\Omega$ отделён от границы: $d=\dist(K,\partial\Omega)>0$.
При $\rho<d$ шар $B(x,\rho)$ для любого $x\in K$ целиком лежит в $\Omega$.}
\label{fig:compact}
\end{figure}
